% Generated by Listas/make_statements.py from pratica-lecture-01.md, item 4. Do not edit.
\textbf{Source.} Muñoz-Garcia \& Toro-Gonzalez (2019), \emph{Strategy and Game Theory: Practice Exercises with Answers}, Ch. 1, Exercise 5 ``Prisoners' Dilemma Game'', p.~8 {[}PDF 21{]}. Reproduced verbatim (statement only; the book's worked answer follows it).



\begin{quote}

Two individuals have been detained for a minor offense and confined in separate cells. The investigators suspect that these individuals are involved in a major crime, and separately offer each prisoner the following deal, as depicted in Fig. 1.8: if you confess while your partner doesn't, you will leave today without serving any time in jail; if you confess and your partner also confesses, you will have to serve 5 years in jail (since prosecutors probably can accumulate more evidence against each prisoner when they both confess); if you don't confess and your partner does, you will have to serve 15 years in jail (since you did not cooperate with the prosecution but your partner's confession provides the police with enough evidence against you); finally, if neither of you confess, you will have to serve one year in jail.



{\def\LTcaptype{none} % do not increment counter

\begin{longtable}[]{@{}lll@{}}

\toprule\noalign{}

Prisoner 1 ~Prisoner 2 & \emph{Confess} & \emph{Not confess} \\

\midrule\noalign{}

\endhead

\bottomrule\noalign{}

\endlastfoot

\emph{Confess} & -5, -5 & 0, -15 \\

\emph{Not confess} & -15, 0 & -1, -1 \\

\end{longtable}

}



\emph{Fig. 1.8} Normal-form of Prisoners' dilemma game



\textbf{Part (a)} Draw the prisoners' dilemma game in its extensive form representation.



\textbf{Part (b)} Mark its initial node, its terminal nodes, and its information set. Why do we represent this information set in the prisoners' dilemma game in its extensive form?



\textbf{Part (c)} How many strategies player 1 has? What about player 2?

\end{quote}

